%=====================================================================
\chapter{First-Order Logic and Resolution}\label{ch:fol}
%=====================================================================
\locator{4}{Part IV --- Knowledge, Logic and Reasoning}

\begin{outcomes}
By the end of this chapter you will be able to:
\begin{tight}
  \item \textbf{Explain} why propositional logic cannot state a general rule, and
        \textbf{quantify} what it costs to try.
  \item \textbf{Express} statements about objects and relations using quantifiers.
  \item \textbf{Unify} two expressions, \textbf{apply} the occurs check, and
        \textbf{say} what goes wrong without it.
  \item \textbf{Convert} a first-order sentence to clause form, including
        skolemisation.
  \item \textbf{Prove} a conclusion by resolution refutation and \textbf{present}
        the proof as an auditable table.
\end{tight}
\end{outcomes}

\begin{prereq}
\chref{propositional} --- entailment and derivation, conjunctive normal form, the
resolution rule, and the equivalence
$\mathit{KB} \models \alpha \Leftrightarrow (\mathit{KB} \wedge \neg\alpha)$
unsatisfiable. Everything here is that machinery with objects added, and
\chref{python} supplies the recursion that unification needs.
\end{prereq}

\begin{keyterms}
first-order logic $\bullet$ constant $\bullet$ variable $\bullet$ function symbol
$\bullet$ term $\bullet$ predicate $\bullet$ atom $\bullet$ universal
quantifier $\bullet$ existential quantifier $\bullet$ substitution $\bullet$
unification $\bullet$ most general unifier $\bullet$ occurs check $\bullet$
standardising apart $\bullet$ skolemisation $\bullet$ clause form $\bullet$
resolution refutation $\bullet$ set of support $\bullet$ subsumption
\end{keyterms}

\section{Where Propositional Logic Runs Out}

\chref{propositional} settled the release gate. Now a rule the delivery team uses
constantly: \emph{if something you depend on is delayed, you are delayed.}

Propositional logic cannot say that. It has no way to talk about ``something'' ---
only about specific, named facts. So the rule has to be written out once per pair:

\begin{quote}
$\mathit{DependsRelease47Payments} \wedge \mathit{DelayedPayments} \Rightarrow
\mathit{DelayedRelease47}$\\
$\mathit{DependsPaymentsAuth} \wedge \mathit{DelayedAuth} \Rightarrow
\mathit{DelayedPayments}$\\
\ldots and so on
\end{quote}

With $n$ services that is $n^{2}$ sentences for one rule, each of which must be
maintained by hand, and a new service means writing $2n$ more. At forty services the
single rule above becomes $1{,}600$ sentences. The knowledge is not large; the
\emph{encoding} is, and it is an encoding no domain expert would agree to maintain.

\dsfig{%
\begin{tikzpicture}[font=\scriptsize,
  svc/.style={rounded corners=3pt, draw=#1, line width=1pt, fill=#1!10,
              text=#1!55!black, font=\scriptsize\bfseries, align=center,
              minimum width=2.0cm, minimum height=0.62cm},
  dep/.style={-{Stealth[length=1.9mm]}, primary!60, line width=1pt},
  el/.style={font=\tiny, text=neutral, inner sep=1.5pt}]

\node[svc=primary] (r47)  at (0,0)    {Release 4.7};
\node[svc=primary] (pay)  at (3.4,0)  {Payments};
\node[svc=accent]  (auth) at (6.8,0)  {Auth};

\draw[dep] (r47) -- node[el, above] {depends on} (pay);
\draw[dep] (pay) -- node[el, above] {depends on} (auth);
\node[font=\scriptsize\bfseries, text=accent, anchor=west] at (7.9,0)
  {delayed};

\node[rounded corners=3pt, fill=lightbg, draw=primary!30, line width=0.9pt,
      text=primary!70!black, font=\scriptsize, text width=6.2cm,
      align=flush left, inner sep=6pt, anchor=north west] at (-0.6,-1.0)
  {\textbf{Propositional} (\chref{propositional}): one sentence per pair.
   $n$ services need $n^{2}$ of them, and adding a service means writing $2n$
   more by hand.};

\node[rounded corners=3pt, fill=greenhl!8, draw=greenhl!55, line width=0.9pt,
      text=greenhl!45!black, font=\scriptsize, text width=6.8cm,
      align=flush left, inner sep=6pt, anchor=north west] at (6.4,-1.0)
  {\textbf{First-order}: one sentence, for every $n$.
   \smallskip
   $\forall x, y \; \mathit{depends}(x,y) \wedge \mathit{delayed}(y)
   \Rightarrow \mathit{delayed}(x)$
   \smallskip
   Adding a service means adding a \emph{fact}, and the rule is untouched.};
\end{tikzpicture}}%
{The delay chain. Quantifiers buy the ability to state a rule once, which is the
difference between a knowledge base a team can maintain and one it
cannot.}{fol-why}

\section{Syntax and Semantics}

First-order logic adds three things to the propositional language.

\begin{tight}
  \item \term{Terms} denote objects: \term{constants} (\texttt{Auth}),
        \term{variables} ($x$), and \term{function symbols} applied to terms
        ($\mathit{owner}(\mathit{Payments})$).
  \item \term{Predicates} denote relations: $\mathit{depends}(x, y)$ is true or
        false of a pair of objects.
  \item \term{Quantifiers}: $\forall x$ (for all) and $\exists x$ (there exists).
\end{tight}

The two quantifiers pair with different connectives, and getting this wrong is the
most common beginner's error:

\[
  \forall x \; \mathit{service}(x) \Rightarrow \mathit{monitored}(x)
  \qquad\text{but}\qquad
  \exists x \; \mathit{service}(x) \wedge \mathit{unmonitored}(x)
\]

Universal with implication; existential with conjunction. Writing
$\forall x\; \mathit{service}(x) \wedge \mathit{monitored}(x)$ claims that
\emph{everything in the universe} is a monitored service --- including Ayesha and
the number seven. Writing
$\exists x\; \mathit{service}(x) \Rightarrow \mathit{unmonitored}(x)$ is satisfied
by any object that is not a service, and so says essentially nothing.

\section{Unification}

Resolution needs to know when two atoms can be made identical. That is
\term{unification}.

\begin{definitionbox}[Unification and the most general unifier]
A \term{substitution} $\theta$ maps variables to terms. $\mathrm{Unify}(p, q)$
returns a substitution with $p\theta = q\theta$, or failure.

Among all unifiers there is a \term{most general} one, which commits to as little as
possible. That matters: a unifier that binds $x$ to \texttt{Auth} when it only
needed to bind $x$ to $y$ throws away conclusions.

The \term{occurs check} refuses to bind a variable to a term containing it.
\end{definitionbox}

\dsfig{%
\begin{tikzpicture}[font=\scriptsize,
  row/.style={rounded corners=3pt, draw=#1, line width=1pt, fill=#1!7,
              text width=13.4cm, align=flush left, inner sep=6pt,
              text=#1!50!black},
  tag/.style={rounded corners=2pt, fill=#1, text=white,
              font=\tiny\bfseries, inner sep=3pt}]

\node[row=greenhl] (r1) at (0,2.2)
  {\\[1pt]\texttt{unify( blocks(x, Auth), \ blocks(Release47, y) )}\\[3pt]
   Match argument by argument: $x$ is a variable, so bind $x
   \mapsto \texttt{Release47}$; then \texttt{Auth} is a constant and $y$ a
   variable, so bind $y \mapsto \texttt{Auth}$.\\[3pt]
   \textbf{Result:} $\theta = \{x/\texttt{Release47},\; y/\texttt{Auth}\}$};

\node[row=accent] (r2) at (0,0.1)
  {\\[1pt]\texttt{unify( blocks(x, x), \ blocks(Auth, Payments) )}\\[3pt]
   Bind $x \mapsto \texttt{Auth}$ from the first argument. The second argument
   then requires \texttt{Auth} to equal \texttt{Payments}, and two distinct
   constants never unify.\\[3pt]
   \textbf{Result: failure.} The same variable twice means \emph{the same
   object} twice.};

\node[row=accent] (r3) at (0,-2.0)
  {\\[1pt]\texttt{unify( x, \ f(x) )}\\[3pt]
   Tempting to bind $x \mapsto f(x)$ --- and that substitution describes no finite
   term. Applying it gives $f(f(f(\ldots)))$, and a prover that accepts it will
   loop or crash.\\[3pt]
   \textbf{Result: failure, by the occurs check.} This is the one case people leave
   out, and it is the one that produces a hang rather than a wrong answer.};

\node[tag=greenhl] at (r1.north west) {SUCCEEDS};
\node[tag=accent]  at (r2.north west) {FAILS};
\node[tag=accent]  at (r3.north west) {FAILS --- OCCURS CHECK};
\end{tikzpicture}}%
{Three unification problems. The third is why every implementation needs an occurs
check, and why many production Prolog systems omit it for speed and document the
risk.}{unify}

\section{Clause Form}

Resolution works on clauses, so first-order sentences must be converted. The
procedure extends \chref{propositional}'s with two steps for quantifiers.

\begin{enumerate}[leftmargin=2.2em, itemsep=2pt, topsep=4pt]
  \item Eliminate $\Leftrightarrow$ and $\Rightarrow$ as before.
  \item Drive $\neg$ inwards. Over quantifiers:
        $\neg \forall x\, P \equiv \exists x\, \neg P$ and
        $\neg \exists x\, P \equiv \forall x\, \neg P$.
  \item \term{Standardise apart}: rename so no two quantifiers bind the same
        variable name.
  \item \term{Skolemise}: replace each existentially quantified variable by a new
        constant, or --- if it lies inside any universal quantifiers --- by a new
        \emph{function} of those universally quantified variables.
  \item Drop the universal quantifiers (everything left is implicitly universal)
        and distribute $\vee$ over $\wedge$.
\end{enumerate}

\begin{examplebox}[Skolemisation, and why the function is necessary]
``Every service has an owner'' is
$\forall x \; \mathit{service}(x) \Rightarrow \exists y \; \mathit{owns}(y, x)$.

\smallskip
Skolemising $y$ to a \emph{constant} \texttt{Sk} gives
$\neg\mathit{service}(x) \vee \mathit{owns}(\texttt{Sk}, x)$, which claims one
person owns every service. That is a different and much stronger statement, and it
is wrong.

\smallskip
$y$ depends on $x$, so it must become a \emph{function} of $x$:
\[
  \neg \mathit{service}(x) \vee \mathit{owns}(\mathit{owner}(x),\, x)
\]
where $\mathit{owner}$ is a new function symbol meaning ``the owner of''. Now each
service has its own. The rule is mechanical: \emph{skolemise to a function of every
universal quantifier you are inside}.
\end{examplebox}

The delay rule needs no skolemisation --- it is purely universal:
\[
  \forall x, y \;\; \mathit{depends}(x,y) \wedge \mathit{delayed}(y)
  \Rightarrow \mathit{delayed}(x)
\]
becomes the single clause
$\neg \mathit{depends}(x,y) \vee \neg \mathit{delayed}(y) \vee
\mathit{delayed}(x)$.

\section{Resolution Refutation}

The rule is \chref{propositional}'s, with unification doing the matching:
\[
  \frac{\alpha \vee \ell \qquad \beta \vee \neg m
        \qquad \theta = \mathrm{Unify}(\ell, m)}
       {(\alpha \vee \beta)\theta}
\]
Before each resolution the two clauses are \term{standardised apart}, so that an $x$
in one is not accidentally identified with an $x$ in the other.

\begin{examplebox}[The knowledge base]
\begin{tight}
  \item R1 \quad $\neg \mathit{depends}(x,y) \vee \neg \mathit{delayed}(y)
        \vee \mathit{delayed}(x)$
  \item F1 \quad $\mathit{depends}(\texttt{Release47}, \texttt{Payments})$
  \item F2 \quad $\mathit{depends}(\texttt{Payments}, \texttt{Auth})$
  \item F3 \quad $\mathit{delayed}(\texttt{Auth})$
\end{tight}
Query: $\mathit{delayed}(\texttt{Release47})$ --- is the release late?
\end{examplebox}

\smallskip
\rowcolors{2}{white}{lightbg}
\noindent\begin{longtable}{@{}L{0.9cm}L{7.8cm}L{6.0cm}@{}}
\toprule
\rowcolor{primary!15}
\textbf{\#} & \textbf{Clause} & \textbf{Derived from} \\
\midrule
\endhead
1 & $\neg \mathit{depends}(x,y) \vee \neg \mathit{delayed}(y) \vee
\mathit{delayed}(x)$ & given (R1) \\
2 & $\mathit{depends}(\texttt{Release47}, \texttt{Payments})$ & given (F1) \\
3 & $\mathit{depends}(\texttt{Payments}, \texttt{Auth})$ & given (F2) \\
4 & $\mathit{delayed}(\texttt{Auth})$ & given (F3) \\
5 & $\neg \mathit{delayed}(\texttt{Release47})$ & \textbf{assumed}: negated query \\
6 & $\neg \mathit{depends}(x, \texttt{Auth}) \vee \mathit{delayed}(x)$ &
1, 4 with $\theta = \{y/\texttt{Auth}\}$ \\
7 & $\mathit{delayed}(\texttt{Payments})$ &
3, 6 with $\theta = \{x/\texttt{Payments}\}$ \\
8 & $\neg \mathit{depends}(x, \texttt{Payments}) \vee \mathit{delayed}(x)$ &
1, 7 with $\theta = \{y/\texttt{Payments}\}$ \\
9 & $\mathit{delayed}(\texttt{Release47})$ &
2, 8 with $\theta = \{x/\texttt{Release47}\}$ \\
10 & $\square$ \quad (the empty clause) & 5, 9 \\
\bottomrule
\caption{A resolution refutation in five derived steps. Each line names its parents
and its unifier, so the proof can be checked by hand --- which is what makes it an
audit artefact rather than an assertion.}\label{tab:fol-proof}
\end{longtable}

\begin{alertbox}[Finding that proof is much harder than checking it]
Table~\ref{tab:fol-proof} is five derived lines. A naive implementation ---
resolve every pair, keep everything new, repeat --- reaches the empty clause in
\textbf{705 resolution steps}.

The standard first remedy is \term{set of support}: require every resolution to
involve a clause descended from the negated query. Nothing else is relevant to
proving \emph{this} conclusion, and the restriction preserves refutation
completeness. It cuts the search to \textbf{202 steps}, a factor of $3.5$.

Even that is far from five, and the reason is visible in the lab's trace: the same
clause is derived repeatedly with different variable names, because standardising
apart creates fresh names each time and a set of clauses does not recognise
alpha-variants as equal. Real provers add \term{subsumption} --- discard a clause
if a more general one is already present --- along with unit preference and
indexing. That is the subject of automated theorem proving, and it is where this
course stops.
\end{alertbox}

\section{What This Buys, and What It Costs}

First-order logic is strictly more expressive than propositional logic, and the
price is paid in decidability.

Propositional entailment is \emph{decidable}: model checking always terminates, if
slowly. First-order entailment is only \emph{semi-decidable}. If the conclusion
follows, resolution will eventually find a proof. If it does not, the procedure may
run forever without ever being able to report that fact --- there are infinitely
many models and it cannot enumerate them to find a countermodel.

So \chref{propositional}'s most practically useful output, the countermodel, is in
general unavailable here. That is a real loss, and it is one reason production
systems often restrict themselves to a decidable fragment --- Horn clauses, which
is what \chref{expert} uses and what makes Prolog work.

\begin{worked}[Following the unifiers through the proof]
The point of this trace is that the general rule R1 is used \emph{three times} with
three different substitutions, which is exactly what propositional logic could not
do.

\smallskip
\textbf{Step 6: resolve R1 with F3.} R1 contains
$\neg \mathit{delayed}(y)$ and F3 is $\mathit{delayed}(\texttt{Auth})$. Unifying
$\mathit{delayed}(y)$ with $\mathit{delayed}(\texttt{Auth})$ gives
$\theta = \{y/\texttt{Auth}\}$. Remove the complementary pair and apply $\theta$
to what remains:
\[
  \neg \mathit{depends}(x, \texttt{Auth}) \vee \mathit{delayed}(x)
\]
Read it in English: \emph{anything that depends on Auth is delayed}. A new general
rule, derived from a general rule and one fact.

\smallskip
\textbf{Step 7: resolve that with F2.} $\mathit{depends}(\texttt{Payments},
\texttt{Auth})$ unifies with $\mathit{depends}(x, \texttt{Auth})$ under
$\theta = \{x/\texttt{Payments}\}$, leaving
$\mathit{delayed}(\texttt{Payments})$. The chain has advanced one link.

\smallskip
\textbf{Step 8: resolve R1 with that.} A \emph{fresh} copy of R1 --- standardised
apart, so its $x$ and $y$ are new variables --- resolves against
$\mathit{delayed}(\texttt{Payments})$ under $\theta = \{y/\texttt{Payments}\}$,
giving $\neg \mathit{depends}(x, \texttt{Payments}) \vee \mathit{delayed}(x)$.

\smallskip
\textbf{Step 9: resolve with F1}, $\theta = \{x/\texttt{Release47}\}$, giving
$\mathit{delayed}(\texttt{Release47})$.

\smallskip
\textbf{Step 10:} that contradicts the assumed $\neg
\mathit{delayed}(\texttt{Release47})$, so the empty clause follows and the query is
proved.

\smallskip
\textbf{Why standardising apart matters.} At step~8, R1's variables must be
renamed. If the $x$ in the fresh copy were identified with the $x$ already bound to
\texttt{Payments} at step~7, the copy would be unusable for
\texttt{Release47} and the proof would fail. The variables in a clause are local to
it, and forgetting that is a bug whose symptom is proofs that mysteriously do not
close.
\end{worked}

\begin{pitfall}
\begin{tight}
  \item \textbf{$\forall$ with $\wedge$.}
        $\forall x\, \mathit{service}(x) \wedge \mathit{monitored}(x)$ claims
        everything in the universe is a monitored service. Use $\Rightarrow$ with
        $\forall$ and $\wedge$ with $\exists$.
  \item \textbf{Omitting the occurs check.} $\mathrm{Unify}(x, f(x))$ then
        succeeds with an infinite term, and the symptom is a hang rather than a
        wrong answer.
  \item \textbf{Skolemising to a constant inside a universal.} ``Every service has
        an owner'' becomes ``one person owns every service''.
  \item \textbf{Forgetting to standardise apart.} Two clauses sharing a variable
        name become accidentally linked, and proofs fail for reasons invisible in
        the output.
  \item \textbf{Expecting a countermodel when the proof fails.} First-order
        entailment is only semi-decidable: failing to find a proof is not evidence
        that none exists.
  \item \textbf{Assuming naive resolution will scale.} $705$ steps for a five-line
        proof. Set of support, subsumption and indexing are not optional extras.
\end{tight}
\end{pitfall}

%---------------------------------------------------------------------
\begin{lab}[Unification and a Resolution Proof]
A real unifier with the occurs check, first-order resolution, and the set-of-support
strategy measured against naive saturation.
\begin{lstlisting}[language=Python]
"""Chapter 12 lab -- first-order resolution on a delay chain.

Representation: a constant is a string; a variable is ("?", name);
a compound term is (functor, arg, ...); a literal is
(sign, predicate, args).
"""


def var(n):
    return ("?", n)


def is_var(t):
    return isinstance(t, tuple) and len(t) == 2 and t[0] == "?"


def is_compound(t):
    return isinstance(t, tuple) and not is_var(t)


def walk(t, s):
    while is_var(t) and t in s:
        t = s[t]
    return t


def occurs(v, t, s):
    """The occurs check: is v inside t? Omit this and unify(x, f(x))
    succeeds with a term that does not exist."""
    t = walk(t, s)
    if t == v:
        return True
    if is_compound(t):
        return any(occurs(v, a, s) for a in t[1:])
    return False


def unify(x, y, s):
    """Most general unifier, or None."""
    if s is None:
        return None
    x, y = walk(x, s), walk(y, s)
    if x == y:
        return s
    if is_var(x):
        if occurs(x, y, s):
            return None
        s = dict(s)
        s[x] = y
        return s
    if is_var(y):
        return unify(y, x, s)
    if is_compound(x) and is_compound(y):
        if x[0] != y[0] or len(x) != len(y):
            return None
        for a, b in zip(x[1:], y[1:]):
            s = unify(a, b, s)
            if s is None:
                return None
        return s
    return None


def subst(t, s):
    t = walk(t, s)
    if is_compound(t):
        return (t[0],) + tuple(subst(a, s) for a in t[1:])
    return t


def lit_subst(lit, s):
    sign, pred, args = lit
    return (sign, pred, tuple(subst(a, s) for a in args))


def standardise(clause, n):
    """Rename every variable, so two clauses never share one."""
    out = []
    for (sign, pred, args) in clause:
        out.append((sign, pred,
                    tuple(("?", "%s#%d" % (a[1], n)) if is_var(a) else a
                          for a in args)))
    return frozenset(out)


def resolve(ci, cj, n):
    ci, cj = standardise(ci, n), standardise(cj, n + 500000)
    out = []
    for li in ci:
        for lj in cj:
            if li[0] == lj[0] or li[1] != lj[1]:
                continue
            if len(li[2]) != len(lj[2]):
                continue
            s, ok = {}, True
            for a, b in zip(li[2], lj[2]):
                s = unify(a, b, s)
                if s is None:
                    ok = False
                    break
            if ok:
                out.append(frozenset(
                    lit_subst(l, s) for l in (ci - {li}) | (cj - {lj})))
    return out


def show_term(a):
    if is_var(a):
        return a[1]
    if is_compound(a):
        return "%s(%s)" % (a[0], ",".join(show_term(x) for x in a[1:]))
    return a


def show(clause):
    if not clause:
        return "[]  (the empty clause)"
    parts = []
    for (sign, pred, args) in clause:
        parts.append(("" if sign else "-") +
                     "%s(%s)" % (pred,
                                 ",".join(show_term(a) for a in args)))
    return " OR ".join(sorted(parts))


# --- the knowledge base -------------------------------------------
x, y = var("x"), var("y")
R1 = frozenset({(False, "depends", (x, y)),
                (False, "delayed", (y,)),
                (True, "delayed", (x,))})
F1 = frozenset({(True, "depends", ("Release47", "Payments"))})
F2 = frozenset({(True, "depends", ("Payments", "Auth"))})
F3 = frozenset({(True, "delayed", ("Auth",))})
KB = [R1, F1, F2, F3]
QUERY = (True, "delayed", ("Release47",))


def refute_naive(kb, query, rounds=6):
    """Resolve every pair, keep everything new, repeat."""
    clauses = list(kb) + [frozenset({(not query[0], query[1], query[2])})]
    steps = 0
    for _ in range(rounds):
        new = []
        for i in range(len(clauses)):
            for j in range(len(clauses)):
                if i == j:
                    continue
                for r in resolve(clauses[i], clauses[j], steps):
                    steps += 1
                    if not r:
                        return True, steps
                    if r not in clauses and r not in new:
                        new.append(r)
        if not new:
            return False, steps
        clauses += new
    return False, steps


def refute_sos(kb, query, rounds=8):
    """Set of support: every resolution must involve a clause
    descended from the negated query. Nothing else is relevant to
    proving THIS conclusion, and completeness is preserved."""
    neg_q = frozenset({(not query[0], query[1], query[2])})
    base, sos = list(kb), [neg_q]
    seen = set(base) | {neg_q}
    steps = 0
    for _ in range(rounds):
        new = []
        for s in sos:
            for b in base + sos:
                if b is s:
                    continue
                for r in resolve(s, b, steps):
                    steps += 1
                    if not r:
                        return True, steps
                    if r not in seen:
                        seen.add(r)
                        new.append(r)
        if not new:
            return False, steps
        sos += new
    return False, steps


if __name__ == "__main__":
    print("=== unification ===")
    cases = [
        ("blocks(x, Auth) vs blocks(Release47, y)",
         ("blocks", x, "Auth"), ("blocks", "Release47", y)),
        ("blocks(x, x)   vs blocks(Auth, Payments)",
         ("blocks", x, x), ("blocks", "Auth", "Payments")),
        ("x              vs f(x)   [occurs check]",
         x, ("f", x)),
    ]
    for label, a, b in cases:
        s = unify(a, b, {})
        if s is None:
            print("   %-42s FAIL" % label)
        else:
            binding = ", ".join("%s/%s" % (k[1], show_term(v))
                                for k, v in sorted(s.items()))
            print("   %-42s {%s}" % (label, binding))

    # the three results the chapter claims
    assert unify(("blocks", x, "Auth"),
                 ("blocks", "Release47", y), {}) is not None
    assert unify(("blocks", x, x),
                 ("blocks", "Auth", "Payments"), {}) is None
    assert unify(x, ("f", x), {}) is None

    print()
    print("=== the proof, by hand ===")
    # step 6: R1 with F3
    s6 = unify(("delayed", y), ("delayed", "Auth"), {})
    c6 = frozenset(lit_subst(l, s6)
                   for l in R1 - {(False, "delayed", (y,))})
    print("   6.", show(c6), "   theta = {y/Auth}")
    # step 7: that with F2
    got = resolve(c6, F2, 1)
    c7 = min(got, key=len)
    print("   7.", show(c7), "   theta = {x/Payments}")
    assert c7 == frozenset({(True, "delayed", ("Payments",))})
    # step 8: R1 with step 7
    c8 = min((c for c in resolve(R1, c7, 2) if len(c) == 2), key=len)
    print("   8.", show(c8), "   theta = {y/Payments}")
    # step 9: that with F1
    c9 = min(resolve(c8, F1, 3), key=len)
    print("   9.", show(c9), "   theta = {x/Release47}")
    assert c9 == frozenset({(True, "delayed", ("Release47",))})
    print("  10. [] -- contradicts the assumed -delayed(Release47)")
    print()
    print("   R1 was used three times, with three different unifiers.")
    print("   That is exactly what propositional logic cannot do.")

    print()
    print("=== finding it automatically ===")
    ok1, n1 = refute_naive(KB, QUERY)
    ok2, n2 = refute_sos(KB, QUERY)
    print("   naive saturation : proved %s in %4d resolution steps"
          % (ok1, n1))
    print("   set of support   : proved %s in %4d resolution steps"
          % (ok2, n2))
    print("   set of support is %.1fx cheaper" % (n1 / n2))
    assert ok1 and ok2 and n2 < n1
    print()
    print("   Both find a five-line proof. Checking a proof is easy;")
    print("   finding one is the hard part, and real provers add")
    print("   subsumption and indexing on top of this.")

    # and a query that does NOT follow
    print()
    print("=== a query that does not follow ===")
    ok3, n3 = refute_sos(KB, (True, "delayed", ("Billing",)), rounds=4)
    print("   delayed(Billing): proved %s after %d steps" % (ok3, n3))
    assert not ok3
    print("   No proof was found -- but first-order entailment is only")
    print("   semi-decidable, so 'no proof yet' is not the same as")
    print("   'no proof exists'. Chapter 11 could hand back a")
    print("   countermodel here; in general this chapter cannot.")
\end{lstlisting}
Then remove the occurs check from \texttt{unify} --- delete the two lines that call
\texttt{occurs} --- and re-run. The first two cases are unchanged and the third now
``succeeds'' with a binding that denotes no finite term. That is the bug, and the
reason it is worth a name.
\end{lab}

%---------------------------------------------------------------------
\begin{checkpoint}
\begin{tightnum}
  \item Write ``every service has an owner'' in first-order logic and convert it to
        clause form. Explain why the existential variable must become a function of
        $x$ rather than a constant.
  \item In the proof of Table~\ref{tab:fol-proof}, rule R1 is used three times.
        State how many propositional sentences would be needed to say the same
        thing over $n$ services, and what happens when a service is added.
  \item \chref{propositional} could return a countermodel when a conclusion did not
        follow. Explain why this chapter generally cannot, and name the property of
        first-order entailment responsible.
\end{tightnum}
\end{checkpoint}

%---------------------------------------------------------------------
\begin{chaptersummary}
\begin{tight}
  \item First-order logic adds \term{terms} (constants, variables, functions),
        \term{predicates} and \term{quantifiers}. One quantified sentence replaces
        the $n^{2}$ propositional ones a general rule would otherwise need.
  \item Pair $\forall$ with $\Rightarrow$ and $\exists$ with $\wedge$. The other
        combinations say something much stronger, or nothing at all.
  \item \term{Unification} finds a substitution making two atoms identical, and the
        \term{most general} one commits to as little as possible. The \term{occurs
        check} refuses to bind a variable inside a term containing it; omitting it
        produces a hang, not a wrong answer.
  \item Clause form adds \term{standardising apart} and \term{skolemisation} to
        \chref{propositional}'s procedure. An existential inside a universal must
        become a \emph{function} of it, or a rule about every service collapses
        into a claim about one.
  \item \term{Resolution} lifts to first-order logic by unifying the complementary
        literals. The proof is an auditable table: each line names its parents and
        its unifier.
  \item Checking a proof is easy; finding one is not. Naive saturation took $705$
        steps for a five-line proof; \term{set of support} reduced it to $202$, and
        real provers add \term{subsumption} and indexing.
  \item First-order entailment is only \emph{semi-decidable}. Failing to find a
        proof is not evidence that none exists, and the countermodel of
        \chref{propositional} is generally unavailable.
\end{tight}
\end{chaptersummary}

\begin{reviewq}
  \item \textbf{(MCQ)} The occurs check prevents: (a)~two constants unifying
        (b)~binding a variable to a term containing it (c)~skolemisation
        (d)~variable capture between clauses.
  \item \textbf{(MCQ)} $\forall x \; \mathit{service}(x) \wedge
        \mathit{monitored}(x)$ claims that: (a)~every service is monitored
        (b)~everything is a monitored service (c)~some service is monitored
        (d)~nothing is monitored.
  \item \textbf{(MCQ)} An existential variable inside a universal quantifier
        skolemises to: (a)~a new constant (b)~a new function of the universally
        quantified variables (c)~the same variable (d)~a predicate.
  \item \textbf{(MCQ)} First-order entailment is: (a)~decidable
        (b)~semi-decidable (c)~undecidable in both directions (d)~decidable only
        for finite domains.
  \item \textbf{(Short)} Explain why clauses must be standardised apart before
        resolution, and describe the symptom when they are not.
  \item \textbf{(Short)} Give the most general unifier of
        $\mathit{owns}(x, \mathit{svc}(y))$ and
        $\mathit{owns}(\texttt{Ayesha}, \mathit{svc}(\texttt{Auth}))$, and say what
        a less general unifier would throw away.
  \item \textbf{(Short)} Set of support cut the search from $705$ steps to $202$.
        State the restriction it imposes and why it does not lose completeness.
  \item \textbf{(Applied)} Add the fact $\mathit{depends}(\texttt{Billing},
        \texttt{Payments})$ to the knowledge base and prove
        $\mathit{delayed}(\texttt{Billing})$ by resolution. Present the proof as a
        table with parents and unifiers, and say which earlier derived clause you
        were able to reuse.
\end{reviewq}
